
\section{问题分析}

从实际问题到模型建立是一种从具体到抽象的思维过程,问题分析这一部分就是沟通这一过程的桥梁,因为它反映了建模者对于问题的认识程度如何,也体现了解决问题的雏形,起着承上启下的作用,也体现出建模者的综合水平。

% 这部分的内容应包括:题目中包含的信息和条件,利用信息和条件对题目做整体分析,确定用什么方法建立模型,一般是每个问题单独分析一小节,分析过程要简明扼要,不需要放结论。
% 建议在文字说明的同时用图形或图表(例如流程图)列出思维过程,这会使你的思维显得很清晰,让人觉得一目了然。

\subsection{问题一的分析}

【问题一是【问题类型:优化/评价/预测/分类/机理】问题,要求【具体求解目标】。从题目条件来看,【关键信息1】和【关键信息2】是核心约束,需要【建模思路概述】。本节将围绕【主要方法】展开,辅以【对比方法】进行交叉验证。】

\subsection{问题二的分析}

【问题二在问题一的基础上进一步考虑【新增约束/变量】,要求【具体求解目标】。需要【说明该问题的核心难点】,并采用【主要方法】解决。XXXXXXXXXXXXXXXXXXXX(一般写 10 行左右即可)】

\subsection{问题三的分析}

【问题三要求【具体求解目标】。XXXXXXXXXXXXXXXXXXXX(一般写 10 行左右即可)】

\subsection{问题四的分析}

【问题四是【建议/综合/扩展】型问题,基于前三问的结果,从【角度1】、【角度2】等维度给出【数量】条【建议/结论】。XXXXXXXXXXXXXXXXXXXX(一般写 10 行左右即可)】

针对以上问题,本文采用了以下整体思路框架进行系统性研究和分析:

\begin{figure}[ht]
  \begin{center}
    \begin{tikzpicture}[x=1cm, y=0.7cm]
      \node[box] (n11) at (0, 0) {问题重述};
      \node[box] (n12) at (5, 0) {问题分析};
      \node[box] (n13) at (10, 0) {模型假设};
      \node[box] (n21) at (0, -3) {数据预处理};
      \node[box] (n22) at (5, -3) {模型建立};
      \node[box] (n23) at (10, -3) {模型求解};
      \node[box] (n31) at (0, -6) {模型检验};
      \node[box] (n32) at (5, -6) {模型评价};
      \node[box] (n33) at (10, -6) {模型改进};
      \draw[arrow] (n11) -- (n12);
      \draw[arrow] (n12) -- (n13);
      \draw[arrow] (n13) -- ++(0,-1.0) -| (n21);
      \draw[arrow] (n21) -- (n22);
      \draw[arrow] (n22) -- (n23);
      \draw[arrow] (n23) -- ++(0,-1.0) -| (n31);
      \draw[arrow] (n31) -- (n32);
      \draw[arrow] (n32) -- (n33);
    \end{tikzpicture}
    \caption{问题的总分析}
    \label{fig:overall_analysis}
  \end{center}
\end{figure}
